% Source-aligned English for printed p.217.
\clearpage
\SourcePage{217}

\LayerTarget{tohoku:proof:lemma-5.6.2}
\LayerTarget{tohoku:construction:equivariant-injective-envelope}
\noindent into a module $M(A)$, which is a (nonadditive!) functor of $A$, the
homomorphism $A\longrightarrow M_1(A)$ being functorial; the procedure is
iterated transfinitely, and one passes to the inductive limit, thereby
obtaining an injective object $M(A)$ which is again a functor in $A$, and an
injection $A\longrightarrow M(A)$ which is a functorial homomorphism. This
construction is well defined once a generator $U$ of $C$ and a suitable
cardinal have been chosen. Moreover, replacing $U$ if necessary by the direct
sum of its transforms under $G$ (compare the proof of Prop. 5.1.1), one may
suppose that $U$ comes from an object of $C(G)$. It is then clear, for evident
reasons of ``transport of structure'', that if $A$ is an object of $C(G)$,
$G$ also acts on $M(A)$, which is thus actually an object of $C(G)$, and the
injection $A\longrightarrow M(A)$ is compatible with the operations of $G$.
This is the desired injection of $A$ into an object of $C(G)$ which is
injective in $C$.

\LayerTarget{tohoku:conclusion:lemma-5.6.2}
This completes the proof; a reader who still has doubts may spell out the
proof in the case $C=C^O$ considered in the lemma, restricting, if preferred,
to the ``purely algebraic'' case in which $X$ consists of a single point.

\LayerTarget{tohoku:transition:main-equivariant-ext-result}
We are now in a position to state the principal result of this section, proved
by the preceding considerations:

\LayerTarget{tohoku:theorem:5.6.3}
\noindent\textsc{Theorem 5.6.3.}\quad \textit{Let $X$ be a space endowed with
a group $G$ of homeomorphisms and a $G$-sheaf of rings $O$. Let $A$ be a
$G$-$O$-module. On the category $C^{O(G)}$ of $G$-$O$-modules there exist
three cohomological spectral functors, converging to the graded functor
$\bigl(\operatorname{Ext}_{O,G}^n(X;A,B)\bigr)$, whose initial terms are,
respectively:}

\LayerTarget{tohoku:equation:5.6.11}
\begin{equation}
\begin{aligned}
\mathrm I_2^{p,q}(B)
  &=H^p\bigl(Y,\mathbf{Ext}_{O,G}^q(A,B)\bigr),\\
\mathrm {II}_2^{p,q}(B)
  &=H^p\bigl(G,\operatorname{Ext}_O^q(X;A,B)\bigr),\\
\mathrm {III}_2^{p,q}(B)
  &=H^p\bigl(X;G,\mathbf{Ext}_O^q(A,B)\bigr).
\end{aligned}
\tag{5.6.11}
\end{equation}

\LayerTarget{tohoku:specialization:theorem-5.6.3}
If $O=A$, the first two spectral sequences reduce to the spectral sequences
of Theorem 5.2.1, while the third spectral sequence is new. If $G$ reduces to
the identity element, then spectral sequences I and III are identical, and
both agree with the spectral sequence of Theorem 4.2.1, whereas spectral
sequence II is trivial. If $X$ consists of a single point, spectral sequences
II and III are identical and, moreover, are nontrivial in general (and,
apparently, useful), whereas spectral sequence I is
trivial.\textsuperscript{11)}

\LayerTarget{tohoku:consequence:edge-homomorphisms-five-term}
The preceding spectral sequences yield edge homomorphisms and five-term exact
sequences, which we shall not write down.

\LayerTarget{tohoku:transition:degeneration-case}
We shall make explicit only one particular case of degeneration:

\LayerTarget{tohoku:footnote:11}
\noindent\textsuperscript{11)} The nontrivial spectral sequence obtained
converges to $\operatorname{Ext}_U(A,B)$, and its initial term is
$H^p\bigl(G,\operatorname{Ext}_O^q(A,B)\bigr)$ ($U$ being the ring generated
by $\mathbf Z(G)$ and the $G$-ring $O$, subject to the commutation relations
$g\lambda g^{-1}=\lambda^g$ for $\lambda\in O$, $g\in G$). Note that this
spectral sequence can be obtained more easily directly, using projective
resolutions of $A$ instead of injective resolutions of $B$. Indeed, it is very
easy to see that if $A$ is a free $U$-module, then $A$ is a free $O$-module
and $\operatorname{Hom}_O(A,B)$ is a $\Gamma^G$-acyclic $G$-module.

\TohokuContinuation{218}
